% \input{\pathsections "sec problem"}

\section{Problem}
\begin{frame}
\frametitle{Demonstration Problem}
\begin{center}
	Let's pose and solve a classic \bl{least squares} problem
\end{center}
\end{frame}

\begin{frame}\frametitle{Validation Tiers of a System Hierarchy}
	\input{\pGlobalLibs figure-hierarchy}
\end{frame}
%
%%  %%  %%  %%  %%
\subsection{Bevington Archetype}
%
\begin{frame}\frametitle{Canonical Reference: Bevington, \rd{2e}}
\center
	Data Reduction and Error Analysis for the Physical Sciences, \rd{2e} \\
	Philip Bevington \\
	\href{https://isbnsearch.org/isbn/9780070051355}{\tiny{ISBN 10: 0070051356  ISBN 13: 9780070051355}} \\[5pt]
	
\includegraphics[ width = 1.25in ]{\pLocalGraphics/bev-cover-zoom} \\
\end{frame}
%
%
\begin{frame}\frametitle{Bevington Table 6.1}
\center
		\includegraphics[ width = 2.0in ]{\pLocalGraphics/bev-data.pdf} \\
\end{frame}
%
%
\begin{frame}\frametitle{Bevington Table 6.1}
\center
	\includegraphics[ width = 2.0in ]{\pLocalGraphics/bev-data-mark-2} \\
\end{frame}
%
\begin{frame}
\frametitle{What we have and what we want}
\begin{table}[htbp]
%\caption{Use data to find best fit to a linear function.}
	\begin{center}
		\begin{tabular}{ccc}
			%
			input & $\Rightarrow$ & output \\\hline
			%
			\includegraphics[ width = 2.0in ]{\pLocalGraphics bev-have} & \raisebox{1.8 cm}{$\Rightarrow$} &
			%
			\includegraphics[ width = 2.0in ]{\pLocalGraphics bev-want} \\
			%
			\{position, temperature\}$_{k=1}^{m}$ & $\Rightarrow$ & intercept and slope \\
			%
			$\lst{x_{k},T_{k}}_{k=1}^{m}$ 
			& $\Rightarrow$ & $\bl{\mat{c}{ a_{0} \\ a_{1} }} \pm \bl{\mat{c}{ \sigma_{0} \\ \sigma_{1}} }$\\
			%
		\end{tabular}
	\end{center}
\label{tab:bevington input-output}
\end{table}%
\end{frame}
%
\begin{frame}\frametitle{The Method of Least Squares.}
	%
\Tmls
\begin{itemize}
	\item \bl{extends toolkit} of linear algebra to inconsistent systems
	\item provides \bl{quantification of the error} in the fit parameters
\end{itemize}
	%
\end{frame}
%

%%  %%  %%  %%  %%
\subsection{Least Squares Exemplar Problem}
%
\begin{frame}\frametitle{\lsq \ Exemplar Problem}
	%
\begin{equation}
		%
	\exemplartall \littlea = \mat{c}{y_{1} \\ y_{2} \\ y_{3}}
		%
\label{eq:exemplar}
\end{equation}
	%
\end{frame}
%
\begin{frame}\frametitle{\lsq \ Exemplar Problem}
	%
Problem:
\begin{equation}
		%
	\mat{cc}{1 & 0 \\ 0 & 1 \\ \rzero & \rzero} \littlea = \mat{c}{y_{1} \\ y_{2} \\ \rd{y_{3}}}
		%
\tag{\ref{eq:exemplar}}
\end{equation}
	%
\end{frame}
%
\begin{frame}\frametitle{\lsq \ Exemplar Solution}
	%
Problem:
\begin{equation}
		%
	\mat{cc}{1 & 0 \\ 0 & 1 \\ \rzero & \rzero} \littlea = \mat{c}{y_{1} \\ y_{2} \\ \rd{y_{3}}}
		%
\tag{\ref{eq:exemplar}}
\end{equation}
	%
Solution:
	%
\begin{equation}
		%
	\bl{\littlea} = \bl{\mat{c}{y_{1} \\ y_{2}}}
		%
\label{eq:solution}
\end{equation}
\end{frame}

\begin{frame}\frametitle{\lsq \ Exemplar Error}
	%
Problem:
\begin{equation}
		%
	\mat{cc}{1 & 0 \\ 0 & 1 \\ \rzero & \rzero} \littlea = \mat{c}{y_{1} \\ y_{2} \\ 0}
		%
\tag{\ref{eq:exemplar}}
\end{equation}
	%
Solution:
	%
\begin{equation}
		%
	\bl{\littlea} = \bl{\mat{c}{y_{1} \\ y_{2}}}
		%
\tag{\ref{eq:solution}}
\end{equation}
%
Residual error:
	%
\begin{equation}
		%
	\rd{r^{2}} = \normts{\bl{\mat{c}{y_{1} \\ y_{2} \\ 0}} - \mat{c}{y_{1} \\ y_{2} \\ \rd{y_{3}}}} = \abs{y_{3}}^{2}
		%
\label{eq:error}
\end{equation}
%
\end{frame}
%
%%  %%  %%  %%  %%
\subsection{\ftola}

\begin{frame}\frametitle{\ftolathm}
	%
\def\mylambda{3.75}
\def\mymid{\mylambda / 2}
\def\vertspace{0.6}
\def\pttop{(0, \mylambda)}
\def\ptright{(\mylambda, 0)}
%
\newcommand{ \axisblueRow }[0]	{\draw[line width=1.25pt,blue,->,>=stealth]	(-0.1,0) -- \ptright	node[above left=0cm and 0cm of \ptright]{$\bornga{*}$};}
\newcommand{ \axisredRow }[0]	{\draw[line width=1.25pt,red, ->,>=stealth]	(0,-0.1) -- \pttop	node[below right=0cm and 0cm of \pttop]{$\rnlla{}$};}
%
\newcommand{ \axisblueCol }[0]	{\draw[line width=1.25pt,blue,->,>=stealth]	(-0.1,0) -- \ptright	node[above left=0cm and 0cm of \ptright]{$\bornga{}$};}
\newcommand{ \axisredCol }[0]	{\draw[line width=1.25pt,red, ->,>=stealth]	(0,-0.1) -- \pttop	node[below right=0cm and 0cm of \pttop]{$\rnlla{*}$};}
%
\begin{table}[htp]
{\scriptsize{
	\begin{center}
		\begin{tabular}{cccc}
				%
			\multicolumn{4}{c}{\bf{The \ftolathm\ }} \\[5pt]
				%
			& $\A{} \in \mathbb{C}^{^{m \times n}}$ \\%\cline{2-2}
				%
			% ROW SPACE
			\begin{tikzpicture}
				\perps{1/4}
				\node[] at ( \mymid,  \mymid + \vertspace )	{\dg{$ \mat{ccc}{a_{11} & \cdots & a_{1n} }$} };
				\node[] at ( \mymid,  \mymid )				{\dg{$ \vdots$}} ;
				\node[] at ( \mymid,  \mymid - \vertspace )		{\dg{$ \mat{ccc}{a_{m1} & \cdots & a_{mn} }$}} ;
				\axisredRow{}
				\axisblueRow{}
			\end{tikzpicture} &
				%
			% MAPPINGS
			\raisebox{1.5cm}{
			\begin{tikzpicture}
				\node () at (\mymid, \mymid + \vertspace) {$\A{} \colon \mathbb{C}^{^{n}} \mapsto \bornga{}$};
				\node () at (\mymid, \mymid - \vertspace) {$ \bornga{*} \mapsfrom \mathbb{C}^{^{m}} \hspace{-1.5mm} \cocolon \A{*} $};
			\end{tikzpicture} } &&
				%
			% COLUMN SPACE
			\begin{tikzpicture}
				\perps{1/4}
				\node[] at ( \mymid,  \mymid )	{ \dg{ $\mat{c}{a_{11} \\ \vdots \\ a_{m1} } \ \cdots \ \mat{c}{a_{1n} \\ \vdots \\ a_{mn} }$ } };
				\axisredCol{}
				\axisblueCol{}
			\end{tikzpicture} \\[1pt]
				%
				%
			Solution Space & && Data Space \\[1pt]
				%
			$\mathbb{C}^{^{n}} = \codecomprow{*}$ &&& $\mathbb{C}^{^{m}} = \codecompcol{*} $ \\
				%
		\end{tabular}
	\end{center}
	}}
%\caption{A geometric view of the \ftola \ showing the four fundamental subspaces induced by a matrix and the mappings between those subspaces.}
\label{tab:ftola-master}
\end{table}
	%
\end{frame}
%
%
\begin{frame}\frametitle{\lsq \ Untangles The Null Spaces}
	%
\begin{figure}[t]
\centering
	\begin{overpic}[ scale = 0.75 ]{\pLocalGraphics colored-masks-02}
			%
	\end{overpic}
\label{fig:Seeing the subspace decomposition}
\caption{Vectors in \href{https://mathworld.wolfram.com/ColumnSpace.html}{\bl{range spaces} $\left( \brnga{} \right.$ or $\left. \brnga{*}\right)$} are colored in blue. Vectors in \href{https://mathworld.wolfram.com/NullSpace.html}{\rd{null spaces} $\paren{\rnlla{*} \text{ or } \rnlla{}}$} are colored in red. Amalgam vectors with components in both spaces are the traditional black.}
\end{figure}
	%
\end{frame}
%
%
\begin{frame}\frametitle{Null Space Nomenclature.}
%
\begin{table}[htp]
%\caption{Confusion with left- and right-handed null spaces.}
\begin{center}
\begin{tabular}{cc}
		%
	\Nsu & \Nsv \\\hline
		%
		%
	%$\rnlla{*}$ & $\rnlla{}$ \\
		%
	\input{\pLocalComponents left-hand.tex} &
	\input{\pLocalComponents right-hand.tex} \\
			% 
		%
	$\brnga{}^{\perp}$ & $\brnga{*}^{\perp}$  \\
		%
	null($\A{*}$) & null($\A{}$) \\
		%
	ker($\A{*}$) & ker($\A{}$) \\
		%
	cokernel & kernel \\
			% 
\end{tabular}
\end{center}
\label{tab:left-and-right}
\end{table}
	%
\end{frame}

%%  %%  %%  %%  %%
\subsection{Least Squares Problem and Solution}

\begin{frame}\frametitle{Problem Statement}
Consider the linear equation
$$\axeb$$
\begin{enumerate}
	\item $\aicmnr$: design matrix
	\item $b\in\cmplxm$: data vector
	\item $\xicn$: solution vector
\end{enumerate}
\end{frame}
%
\begin{frame}\frametitle{\lsq \ Solution}
The \bl{\lsq \ Solution} is defined as
$$\xlsdefault$$
\end{frame}


\begin{frame}\frametitle{\lsq \ Solution}
The \bl{\lsq \ Solution} is expressed as
\begin{equation}
		\begin{array}{cccccl}
			\xls 
			%& = & \xblue{}} & + & \xred \\
			& = & \bsolnlsbl{b}{y}  \\
			& = & \Apb & + & \pna{}y \\%\myline
			& = & \xp & + & \xh \\
%			& = & \bl{x_{\mathcal{R}}} \\
%			& = & \bl{x_{\scriptstyle\rnga{*}}} \\
			& = & \xblue{} & + & \xred
		\end{array}
\label{eq:definition:general lsq solution}
\end{equation}
\end{frame}

\begin{frame}\frametitle{Dramatis Personae}
%
\lsq \ Residual Error Vector
%
\begin{equation}
	\rd{r(x)} = \Ax - b
\end{equation}
%
\lsq \ Merit Function
%
\begin{equation}
		%
	M(x) = \innerecho{\rd{r(\bk{x})}} = \innerecho{\axmb}
		%
%\label{eq:}
\end{equation}
%
Least Total Squared Error
%
\begin{equation}
		%
	t^{2} = \innerecho{\rd{r(\bl{\xls})}}
		%
%\label{eq:}
\end{equation}
%
\lsq \ Uncertainty \mg{(Standard Deviation)}
%
\begin{equation}
		%
	\sigma^{2}_{k} =  \frac{t^{2}} {m-n} \wxi{*}_{k,k}, \quad \kto{n}
		%
%\label{eq:}
\end{equation}
%
\end{frame}

%%  %%  %%  %%  %%
\subsection{The Fredholm Alternative}
%
\begin{frame}\frametitle{A Vital Theorem}
	%
\input{\pGlobalTheorems "thm Fredholm"}
\pause
	%
	\ \\
Colloquially, \bl{either} $\axmb = 0$ \bl{or} $\axmb > 0$
	%
\end{frame}
%
%
\begin{frame}\frametitle{Fredholm Alternative: Colloquially}
	%
Fredholm's Alternatives are \\
	\bl{either} $\axmb = 0$ \\
	\bl{or} \phantom{mj}\, $\axmb > 0$
	%
\end{frame}

\def\mylambda{2.75}
\def\data{\mylambda * 0.85}
\begin{frame}\frametitle{Fredholm Alternative in Pictures}
%
\begin{table}[t]
%\caption
%[The Fredholm Alternative in pictures]
%{The Fredholm Alternative in pictures. The theorem is statement about the residency of the data vector.}
	\begin{center}
		\begin{tabular}{ccc}
			%
			either && or \\
			%
			$b=\datarange$ && $b=\dataresolved$ \\\hline
			%
			\input{\pLocalComponents "tikz app fredholm-I"} && \input{\pLocalComponents "tikz app fredholm-II"} \\
			%
			$\axeb$ has solution && $\A{*}\rd{y}=0$ and $\rd{y}^{*}b\ne 0$\\
			%
		\end{tabular}
	\end{center}
\label{tab:fredholm-alternatives}
\end{table}%
	%
\end{frame}
%
\begin{frame}\frametitle{Fredholm's Duplex: Residency of the Data Vector}
	%
\begin{figure}[]
\center
\begin{overpic}[ scale = 0.35 ]{\pLocalGraphics fredholm-duplex.png}
\end{overpic}
%\caption
%{The Fredholm duplex. The Fredholm Alternative is statement about the \bl{residency of the data vector}. Data vectors residing entirely in the range space live in apartment I. Data vectors with components in the null space live in apartment II.}
\label{fig:fredholm-duplex}
\end{figure}
	%
\end{frame}
%
\begin{frame}\frametitle{Fredholm's Duplex: Residency of the Data Vector}
	%
\begin{figure}[]
\center
\begin{overpic}[ scale = 0.35 ]{\pLocalGraphics fredholm-duplex.png}
\put(21, -2){\colorbox{white}{$b=\datarange{}$ }}
\put(61, -3){\colorbox{white}{$b=\dataresolved{}$ }}
\end{overpic}
%\caption
%{The Fredholm duplex. The Fredholm Alternative is statement about the \bl{residency of the data vector}. Data vectors residing entirely in the range space live in apartment I. Data vectors with components in the null space live in apartment II.}
\label{fig:fredholm-duplex}
\end{figure}
	%
\end{frame}


\endinput  %  ==  ==  ==  ==  ==  ==  ==  ==  ==
